\begin{exercise}[Why local particles have Bose or Fermi statistics]{I-CH02-001}
Let
\[
  \mathcal C_n(\mathbb R^{D-1})
  =\{(\mathbf x_1,\ldots,\mathbf x_n):\mathbf x_r\ne\mathbf x_s\}/S_n
\]
be the configuration space of $n$ identical particles in $D-1\geq3$ spatial
dimensions. Assume that the particles can be created in bounded, mutually
space-like separated regions and that adiabatic exchanges act unitarily on
their states.

Show that exchange paths give the permutation group $S_n$, and determine all
one-dimensional unitary exchange laws. Relate the two answers to commuting and
anticommuting local creation operators. Higher-dimensional representations of
$S_n$ describe parastatistics in particle quantum mechanics; explain why, under
the usual local-QFT assumptions of positive energy, a unique vacuum, finite
statistics, and transportable localized charges, such a representation is
equivalent to ordinary Bose or Fermi fields with a finite internal index and a
compact gauge redundancy. Finally, identify the topological step that changes
in two spatial dimensions.
\end{exercise}
